% chapters/state_of_art/sota_conclusions.tex
% Sección 2.3: Conclusiones del estado del arte (~0.5-1 página)
% ============================================================================
%
% GUÍA: Sintetizar las conclusiones de la revisión:
%
%   1. Los sistemas existentes (Gradescope, Crowdmark, etc.) son SaaS
%      cerrados que no permiten autoalojamiento, lo que plantea problemas
%      de privacidad (RGPD) y dependencia de terceros.
%
%   2. Las plataformas LMS (Moodle, Canvas) no están diseñadas para
%      corregir exámenes manuscritos con múltiples correctores.
%
%   3. Se justifica el desarrollo de un sistema propio, autoalojable,
%      que combine OCR, corrección colaborativa y trazabilidad completa.
%
%   4. Las tecnologías elegidas (Django, Celery, MinIO, Docker) permiten
%      construir una arquitectura distribuida escalable y mantenible.
%
%   Cerrar con la transición al siguiente capítulo:
%   "A partir de este análisis, el siguiente capítulo recoge los requisitos
%    funcionales y no funcionales que debe satisfacer el backend del \ac{scdee}."
% ============================================================================

% TODO: Redactar las conclusiones del estado del arte.
